% Chapter 4: Words 751 - 1000
\chapter{Intermediate Milestone}
\label{chap:chapter4}

{\small\sffamily\color{mutedtext} Key transition vocabulary reaching the 1,000-word milestone for conversational fluency.\par}

\vspace{0.4em}
\markboth{Chapter 4: Intermediate Milestone}{Words 751--1000}

\begin{multicols}{2}
\vocentry{751}{Роль}{rol’}{A role, a part}{Она получила главную роль в новом фильме.}{She got a starring role in a new film.}
\vocentry{752}{Рост}{rost}{Growth; height}{Рост цен на нефть снизился в этом году.}{This year the growth of oil price has decreased.}
\vocentry{753}{Природа}{pri'roda}{Nature}{Весной природа просыпается после зимнего сна.}{In spring, nature wakes up after the winter sleep.}
\vocentry{754}{Политический}{pali'titcheskij}{Political}{Политический климат в стране ухудшился.}{The political climate in the country has gotten worse.}
\vocentry{755}{Точка}{'tochka}{A point; a dot}{Это событие – стартовая точка конфликта.}{This event is the starting point of the conflict.}
\vocentry{756}{Звезда}{zvez'da}{A star}{Люди говорят, что падающая звезда может осуществить желание.}{People say that a shooting star can make a wish come true.}
\vocentry{757}{Петь}{pet’}{To sing}{Многие люди любят петь в ду́ше.}{Many people like to sing in the shower.}
\vocentry{758}{Садиться}{sa'ditsa}{To sit down}{На эту скамейку нельзя садиться – она окрашена.}{One can’t sit down on this bench – it’s been painted.}
\vocentry{759}{Фамилия}{fa'milija}{A surname, last name}{Я не могу понять, кто он по национальности, а его фамилия ни о чём не говорит.}{I can’t understand what his nationality is and his surname doesn’t say anything.}
\vocentry{760}{Характер}{ha'rachter}{Character, temper}{Она очень красивая девушка, но её характер всё портит.}{She’s a very beautiful girl but her character spoils everything.}
\vocentry{761}{Пожалуйста}{pa'zhalusta}{Please; you are welcome}{Пожалуйста, перестаньте шуметь.}{Please, stop making noise.}
\vocentry{762}{Выше}{'vyshe}{Higher}{Зимой наши расходы выше.}{Our outgoings are higher in winter.}
\vocentry{763}{Офицер}{afi'tser}{An officer (military)}{Офицер повёл себя как настоящий герой.}{The officer behaved like a true hero.}
\vocentry{764}{Толпа}{tol'pa}{A crowd}{Перед посольством собралась толпа людей.}{A crowd of people gathered in front of the embassy.}
\vocentry{765}{Перестать}{pere'stat’}{To stop, to cease}{Я не могу перестать думать о тебе.}{I can’t stop thinking about you.}
\vocentry{766}{Придтись}{prid'tis’}{To fall on (a date) (a variation of ‘прийтись’)}{В этом году мой день рождения должен придтись на вторник.}{My birthday should fall on Tuesday this year.}
\vocentry{767}{Уровень}{'urəven’}{A level}{Уровень образованности в стране повысился.}{The level of education in the country has increased.}
\vocentry{768}{Неизвестный}{neiz'vesnyj}{Unknown}{Космос – это другой, неизвестный мир.}{Space is a different unknown world.}
\vocentry{769}{Кресло}{'kreslə}{An armchair}{Мы купили красивое винтажное кресло.}{We’ve bought a beautiful vintage armchair.}
\vocentry{770}{Баба}{'baba}{Colloquial, sometimes with a rude connotation for ‘woman’, ‘broad’}{Кто эта баба, с которой ты сидел в кафе?}{Who’s that broad you were sitting with in a cafe?}
\vocentry{771}{Секунда}{se'kunda}{A second}{Мне нужна всего лишь секунда вашего времени.}{I need just one second of your time.}
\vocentry{772}{Пожаловать}{pa'zhaləvat’}{Most often used in a set expression ‘Добро пожаловать!’ - Welcome!; to grant}{Дамы и господа, добро пожаловать на наш ежегодный фестиваль!}{Ladies and gentlemen, welcome to our annual festival!}
\vocentry{773}{Банк}{bank}{A bank}{Во время кризиса этот банк потерял много клиентов.}{This bank lost lots of clients during the crisis.}
\vocentry{774}{Опыт}{'opyt}{Experience; an experiment}{Опыт работы в продажах помог ему получить работу.}{The experience of working in sales helped him get the job.}
\vocentry{775}{Тихий}{'tihij}{Quiet, still}{Тихий звук скрипки создал очень приятную атмосферу.}{The quiet sound of the violin created a very pleasant atmosphere.}
\vocentry{776}{Поскольку}{pas'kolku}{Since}{Поскольку эта машина очень дорогая, мы возьмём кредит.}{Since this car is very expensive we’ll take a loan.}
\vocentry{777}{Сапог}{sa'pog}{A boot}{Посмотри, твой левый сапог весь в грязи.}{Look, your left boot is all covered in mud.}
\vocentry{778}{Правило}{'pravilo}{A rule}{Это грамматическое правило очень простое.}{This grammar rule is very simple.}
\vocentry{779}{Стекло}{stek'lo}{Glass}{А вы знали, что стекло делают из песка?}{Did you know that glass is made of sand?}
\vocentry{780}{Получать}{palu'chat’}{To receive, to get}{Я буду получать твою почту, пока ты будешь в командировке.}{I’ll receive your mail while you’re on a business trip.}
\vocentry{781}{Внутренний}{'vnutrennij}{Inner}{Внутренний мир человека – это загадка.}{A person’s inner world is a mystery.}
\vocentry{782}{Дочь}{dotch}{A daughter}{Наша дочь заканчивает школу в этом году.}{Our daughter graduates from school this year.}
\vocentry{783}{Называться}{nazy'vatsa}{To be called; to be named (impersonal)}{Наш ресторан будет называться “Мелодия”.}{Our restaurant is going to be called ‘Melody’.}
\vocentry{784}{Надеяться}{na'dejatsa}{To hope}{Доктор сказал, что мы можем надеяться на её скорое выздоровление.}{The doctor said that we could hope for her quick recovery.}
\vocentry{785}{Член}{tchlen}{A member; colloquial for penis}{Он член влиятельной международной организации.}{He’s a member of an influential international organisation.}
\vocentry{786}{Протянуть}{pratja'nut’}{To reach out; to hold out}{Фокусник попросил меня протянуть руку и закрыть глаза.}{The illusionist asked me to reach out my hand and close my eyes.}
\vocentry{787}{Государственный}{gasu'darsvennyj}{State (adj)}{Этот государственный институт отвечает за поддержку семей с детьми.}{This state institution is responsible for the support of families with children.}
\vocentry{788}{Десяток}{de'sjatək}{A dozen (literally means 10)}{Утро только началось, а я уже получил с десяток деловых звонков.}{The morning has just begun and I’ve already received about a dozen business calls.}
\vocentry{789}{Глубокий}{glu'bokij}{Deep (adj)}{Этот ручей не глубокий, и мы можем легко перейти его.}{This stream isn’t deep and we can easily cross it.}
\vocentry{790}{Цветок}{tsve'tok}{A flower}{На конверте был нарисован нежный розовый цветок.}{A tender pink flower was painted on the envelope.}
\vocentry{791}{Ах}{ah}{Ah!, oh!}{Ах, какой прекрасный букет!}{Ah, what a wonderful bouquet!}
\vocentry{792}{Желание}{zhe'lanije}{A wish, desire}{Загадай желание и никому не рассказывай, а то не сбудется!}{Make a wish and don’t tell anyone or else it won’t come true!}
\vocentry{793}{Дождь}{dozhd’}{Rain}{Дождь был таким сильным, что затопил улицы.}{The rain was so hard that it flooded the streets.}
\vocentry{794}{Впереди}{vpere'di}{Ahead; in front of}{У нас столько счастливых дней впереди!}{We’ve got so many happy days ahead!}
\vocentry{795}{Подходить}{padha'dit’}{To fit, to suit, to match; to come up to}{Этот размер не может тебе подходить – ты гораздо выше.}{This size can’t fit you – you’re much taller.}
\vocentry{796}{Много}{'mnogə}{Much, many, a lot of}{У меня не много времени, так что давайте перейдём к делу.}{I don’t have much time so let’s get down to business.}
\vocentry{797}{Лоб}{lob}{A forehead}{Его лоб был покрыт капельками пота.}{His forehead was covered with drops of sweat.}
\vocentry{798}{Улыбка}{u'lybka}{A smile}{У неё была солнечная улыбка и добрые глаза.}{She had a sunny smile and kind eyes.}
\vocentry{799}{Борьба}{bar’'ba}{Struggle, fight}{Борьба с самим собой – самая сложная борьба.}{The struggle against yourself is the most difficult struggle.}
\vocentry{800}{Ворот}{'vorət}{A collar}{Ворот его рубашки был мокрым и грязным.}{The collar of his shirt was wet and dirty.}
\vocentry{801}{Ящик}{'jastchik}{A box; a drawer}{Мальчик спрятал ящик со своими сокровищами под кроватью.}{The boy hid a box with his treasures under the bed.}
\vocentry{802}{Этаж}{ɛ'tazh}{A floor, a storey}{Первый этаж нашего дома был построен за неделю.}{The first floor of our house was built within a week.}
\vocentry{803}{Служить}{slu'zhit’}{To serve, to serve as}{После школы он собирался служить в армии.}{After school, he was going to serve in the army.}
\vocentry{804}{Вновь}{vnof’}{Again, anew (literary style)}{Мы вновь собрались здесь, чтобы вспомнить великого поэта.}{We gathered here again to recall the great poet.}
\vocentry{805}{Голубой}{galu'boj}{Blue; gay}{Этот голубой галстук не подходит к коричневому ремню.}{This blue tie doesn’t match the brown belt.}
\vocentry{806}{Нечего}{'netchego}{There is nothing to}{Ему нечего делать, и он умирает от скуки.}{There is nothing for him to do and he’s dying of boredom.}
\vocentry{807}{Революция}{reva'lutsija}{A revolution}{Мне кажется, что любая революция – это очень агрессивный метод.}{It seems to me that any revolution is a very aggressive method.}
\vocentry{808}{Впервые}{fper'vyje}{For the first time}{Когда мне было шесть, я впервые увидела настоящий водопад.}{When I was six I saw a real waterfall for the first time.}
\vocentry{809}{Сосед}{sa'sed}{A neighbor}{Мой сосед редко устраивает шумные вечеринки.}{My neighbor rarely makes noisy parties.}
\vocentry{810}{Сестра}{ses'tra}{A sister}{Его сестра вегетарианка.}{His sister is a vegetarian.}
\vocentry{811}{Долгий}{'dolgij}{Long}{У нас был долгий разговор о наших отношениях.}{We had a long talk about our relationship.}
\vocentry{812}{Чей}{tchej}{Whose}{Вы не знаете, чей это зонт?}{Do you know whose umbrella it is?}
\vocentry{813}{Поверить}{pa'verit’}{To believe (perfective)}{Я не могу поверить, что ты выиграла в лотерею.}{I can’t believe that you’ve won the lottery.}
\vocentry{814}{Ситуация}{situ'atsija}{A situation}{Ситуация сложная, но я знаю, что делать.}{The situation is difficult but I know what to do.}
\vocentry{815}{Взглянуть}{vzglja'nut’}{To have a look, to glance}{Детям было любопытно взглянуть на новорождённых щенков.}{The children were curious to have a look at the newborn puppies.}
\vocentry{816}{Слабый}{'slabyj}{Weak, feeble}{Пока мой папа болел, он был такой слабый, что едва мог есть сам.}{When my father was sick he was so weak that he could barely eat himself.}
\vocentry{817}{Количество}{ka'litchestvə}{Number, quantity, amount}{В этом году количество выпускников, проваливших тест, уменьшилось.}{The number of graduates who failed the test decreased this year.}
\vocentry{818}{Вызывать}{vyzy'vat’}{To cause; to call}{Арахис может вызывать аллергию у некоторых людей.}{Peanuts can cause allergy in some people.}
\vocentry{819}{Уверенный}{u'verennyj}{Confident, sure}{Я изменил своё мнение, когда увидел его уверенный взгляд.}{I changed my mind when I saw his confident look.}
\vocentry{820}{Выход}{'vyhəd}{A way out; an exit}{Мы должны подумать и найти выход из этой ситуации.}{We must think and find a way out of this situation.}
\vocentry{821}{Совет}{sa'vet}{Advice}{Его совет помог мне принять правильное решение.}{His advice helped me take the right decision.}
\vocentry{822}{Дурак}{du'rak}{A fool, an idiot}{Ты такой дурак, что упустил эту возможность!}{You are such a fool to have missed this opportunity!}
\vocentry{823}{Любимый}{lju'bimyj}{Favorite}{Рождество – мой любимый праздник.}{Christmas is my favorite holiday.}
\vocentry{824}{Союз}{sa'juz}{A union; an alliance}{Европейский Союз поддерживает беженцев из восточных стран.}{The European Union supports refugees from eastern countries.}
\vocentry{825}{Лето}{'letə}{Summer}{Прошлое лето было дождливым и прохладным.}{The last summer was rainy and cool.}
\vocentry{826}{Ожидать}{azhi'dat’}{To expect, to anticipate}{Вы можете ожидать, что перевод будет готов к среде.}{You can expect the translation to be ready by Wednesday.}
\vocentry{827}{Огород}{aga'rod}{A vegetable garden}{Здорово иметь свой огород! Все овощи натуральные и свежие.}{It’s great to have your own vegetable garden! All the vegetables are natural and fresh.}
\vocentry{828}{Висеть}{vi'set’}{To hang}{Почему твоё платье лежит на полу? Оно должно висеть в шкафу.}{Why is your dress lying on the floor? It must hang in the wardrobe.}
\vocentry{829}{Граница}{gra'nitsa}{A border; a boundary}{Граница между двумя государствами расположена вдоль реки.}{The border between the two states is situated along a river.}
\vocentry{830}{Цвет}{tsvet}{A color}{Цвет глаз может сказать что-нибудь о человеке?}{Can the color of the eyes say anything about the person?}
\vocentry{831}{Серьёзный}{serjeznyj}{Serious, earnest}{Для школьника он очень серьёзный и спокойный.}{He’s very serious and calm for a schoolboy.}
\vocentry{832}{Создать}{saz'dat’}{To create}{С помощью этой программы легко создать новый проект.}{With the help of this program it’s easy to create a new project.}
\vocentry{833}{Интересный}{inte'resnyj}{Interesting}{Это очень интересный роман! Я прочёл его за три дня!}{It’s a very interesting novel! I read it in three days!}
\vocentry{834}{Свобода}{sva'boda}{Freedom}{Свобода слова важна для развития демократического общества.}{Freedom of speech is important for the development of democratic society.}
\vocentry{835}{Зато}{za'to}{On the other hand}{Он не зарабатывает много, но зато много времени проводит с семьёй.}{He doesn’t earn much but on the other hand he spends much time with his family.}
\vocentry{836}{Стул}{stul}{A chair}{Я предложил посетителю стул, но он отказался.}{I offered the visitor a chair but he refused.}
\vocentry{837}{Уехать}{u'jehat’}{To leave, to go away; to drive off (perfective)}{Она решила уехать из страны и начать новую жизнь.}{She decided to leave the country and start a new life.}
\vocentry{838}{Поезд}{'poezd}{A train}{Наш поезд прибывает ровно в пять.}{Our train arrives at five sharp.}
\vocentry{839}{Музыка}{'muzyka}{Music}{Классическая музыка не для меня. Я её не понимаю.}{Classical music is not for me. I don’t understand it.}
\vocentry{840}{Тень}{ten’}{Shade; a shadow}{В солнечный день нелегко найти тень.}{It’s not easy to find shade on a sunny day.}
\vocentry{841}{Лошадь}{'loshat’}{A horse}{Моя внучка всегда хотела иметь лошадь.}{My granddaughter has always wanted to have a horse.}
\vocentry{842}{Поле}{'pole}{A field}{За домом было большое пшеничное поле.}{There was a big field of wheat behind the house.}
\vocentry{843}{Выглядеть}{'vygljadet’}{To look (to have an appearance, not the process of watching)}{Она всегда старается модно выглядеть.}{She always tries to look fashionable.}
\vocentry{844}{Учиться}{u'tchitsa}{To study}{Ты должен усердно учиться, если хочешь поступить в университет.}{You must study hard if you want to enter a university.}
\vocentry{845}{Левый}{'levyj}{Left, left-hand}{Посмотри, твой левый карман порван.}{Look, your left pocket is torn.}
\vocentry{846}{Разговаривать}{razga'varivat’}{To talk, to speak, to converse (with)}{Мы с сестрой можем часами разговаривать друг с другом.}{My sister and I can talk to each other for hours.}
\vocentry{847}{Детский}{'detskij}{Child’s, child (as an adj)}{Детский смех делает жизнь счастливее.}{A child’s laughter makes life happier.}
\vocentry{848}{Тип}{tip}{A type, a kind; slang for ‘dude’, ‘guy’}{Какой тип счёта Вы хотите открыть?}{What type of account do you want to open?}
\vocentry{849}{Суд}{sud}{Court; judgement}{Суд постановил, что он не виновен.}{The court stated that he’s not guilty.}
\vocentry{850}{Связанный}{'svjazannyj}{Connected (with)}{Я нашла файл, связанный с твоим делом.}{I’ve found a file connected with your case.}
\vocentry{851}{Горячий}{gar'jatchij}{Hot}{Будь осторожен – чайник горячий!}{Be careful – the kettle is hot!}
\vocentry{852}{Площадь}{'plostchad’}{A square (open space in a town or a city); an area}{Это историческое место города – площадь Свободы.}{It’s an historical place of the city – Freedom Square.}
\vocentry{853}{Помогать}{pama'gat’}{To help}{Настоящие друзья должны всегда помогать друг другу.}{True friends should always help each other.}
\vocentry{854}{Счастливый}{stchas'livyj}{Happy}{Я помню тот счастливый день, когда родился мой внук.}{I remember that happy day when my grandson was born.}
\vocentry{855}{Повернуться}{paver'nutsa}{To turn around}{Я понял, что мне лучше повернуться и уйти.}{I realized that I’d better turn around and leave.}
\vocentry{856}{Позволить}{paz'volit’}{To afford; to allow, to permit}{В этом году мы не можем позволить себе отпуск за границей.}{This year we can’t afford a holiday abroad.}
\vocentry{857}{Встретить}{'vstretit’}{To meet, to encounter}{Я не ожидал встретить её в этом незнакомом городе.}{I didn’t expect to meet her in that unfamiliar town.}
\vocentry{858}{Радость}{'radost’}{Joy}{Это большая радость – встретиться снова после стольких лет.}{It’s a great joy to meet again after so many years.}
\vocentry{859}{Острый}{'ostryj}{Sharp; acute}{Нож был очень острый, и я порезался.}{The knife was very sharp and I cut myself.}
\vocentry{860}{Возраст}{'vozrast}{Age}{Его возраст совсем не соответствует его внешности.}{His age doesn’t correspond to his appearance at all.}
\vocentry{861}{Орган}{'organ}{An organ (anatomy)}{Печень – это орган, который работает как фильтр.}{The liver is an organ that works as a filter.}
\vocentry{862}{Карта}{'karta}{A map}{Карта была неправильной, и мы заблудились.}{The map was incorrect and we lost our way.}
\vocentry{863}{Входить}{vha'dit’}{To enter, to go into; to be included in}{Нельзя входить в чью-либо комнату, не постучавшись.}{One can’t enter anybody’s room without knocking.}
\vocentry{864}{Обнаружить}{abna'ruzhit’}{To find, to find out}{Я был удивлён обнаружить мобильный сестры в своём рюкзаке.}{I was surprised to find my sister’s mobile in my backpack.}
\vocentry{865}{Король}{ka'rol’}{A king}{Старый король был справедливым и мудрым.}{The old king was just and wise.}
\vocentry{866}{Слава}{'slava}{Fame}{Слава и деньги часто портят людей.}{Fame and money often spoil people.}
\vocentry{867}{Полковник}{pal'kovnik}{A colonel (military)}{Полковник был одиноким человеком без семьи и детей.}{The colonel was a lonely man without a family and kids.}
\vocentry{868}{Мелкий}{'melkij}{Little, small, minute; shallow}{Мелкий дождь и холодная погода ещё больше ухудшили моё настроение.}{The little rain and the cold weather made my mood even worse.}
\vocentry{869}{Бок}{bok}{A side}{Ветеринар попросил меня перевернуть собаку на левый бок.}{The vet asked me to turn the dog on its left side.}
\vocentry{870}{Цена}{tse'na}{Price}{Цена этого костюма вполне разумная.}{The price of this suit is quite reasonable.}
\vocentry{871}{Информация}{infar'matsija}{Information}{Эта информация очень полезна для нашего исследования.}{This information is very useful for our research.}
\vocentry{872}{Мозг}{mosk}{Brain}{Сегодня мой мозг просто отказывается работать.}{My brain simply refuses to work today.}
\vocentry{873}{Удовольствие}{uda'vol’stvije}{A pleasure}{Слушать эту музыку вживую – удовольствие.}{Listening to this music live is a pleasure.}
\vocentry{874}{Воля}{'volja}{Will, willpower; freedom}{Его сильная воля помогла ему добиться своей цели.}{His strong will helped him achieve his goal.}
\vocentry{875}{Область}{'oblast’}{A region, a province; a sphere}{Эта область страны богата лесами и озёрами.}{This region of the country is rich in forests and lakes.}
\vocentry{876}{Крыша}{'krysha}{A roof}{Крыша нашего загородного дома течёт, и мы будем менять её.}{The roof of our country house is leaking and we’re going to fix it.}
\vocentry{877}{Нести}{nes'ti}{To carry}{Она не может сама нести эти сумки. Помоги ей.}{She can’t carry these bags herself. Help her.}
\vocentry{878}{Обратно}{ab'ratnə}{Back}{Мы не довольны качеством и хотим наши деньги обратно.}{We’re not happy with the quality and want our money back.}
\vocentry{879}{Современный}{savre'mennyj}{Modern, contemporary}{Современный мир меняется быстрее, чем век назад.}{The modern world changes faster than a century ago.}
\vocentry{880}{Дама}{'dama}{A lady}{Эта пожилая дама много помогает бездомным.}{This elderly lady helps the homeless a lot.}
\vocentry{881}{Семь}{sem’}{Seven}{Многие считают, что семь – это магическое число.}{Many people consider seven to be a magical number.}
\vocentry{882}{Весёлый}{ves'jolyj}{Merry, cheerful, happy}{Весёлый клоун принёс детям много радости.}{The merry clown brought the children much joy.}
\vocentry{883}{Прислать}{pris'lat’}{To send (to, to this place) (perfective)}{Я обещаю прислать тебе открытку, когда буду на курорте.}{I promise to send you a postcard when I’m at the resort.}
\vocentry{884}{Сад}{sad}{A garden, an orchard}{Папа много работает, чтобы наш сад был таким красивым.}{Dad works a lot for our garden to be so beautiful.}
\vocentry{885}{Правительство}{pra'vitel’stvə}{A government}{Новое правительство менее популярно среди народа.}{The new government is less popular with people.}
\vocentry{886}{Милый}{'milyj}{Cute}{Этот щенок такой милый! Давай оставим его!}{This puppy is so cute! Let's keep it!}
\vocentry{887}{Относиться}{atnə'sitsa}{To treat}{Постарайся относиться к ней с пониманием.}{Try to treat her with understanding.}
\vocentry{888}{Возникать}{vazni'kat’}{To arise, to emerge}{При решении этой задачи могут возникать некоторые вопросы.}{While solving this task, some questions may arise.}
\vocentry{889}{Мол}{mol}{He says, they say (colloquial)}{Они не одолжили мне денег, потому что я, мол, слишком мало зарабатываю.}{They didn’t lend me any money because they say I earn too little.}
\vocentry{890}{Повторить}{pavta'rit’}{To repeat}{Учитель попросил учеников повторить последнее предложение.}{The teacher asked the pupils to repeat the last sentence.}
\vocentry{891}{Название}{naz'vanije}{A name, a title}{Ты помнишь название того блюда, что мы готовили на прошлой неделе?}{Do you remember the name of the dish we cooked last week?}
\vocentry{892}{Средний}{'srednij}{Average; middle}{Каков средний доход по стране в этом месяце?}{What’s the average income across the country this month?}
\vocentry{893}{Пример}{pri'mer}{An example}{Ты можешь привести пример, чтобы я лучше понял правило?}{Can you give an example so that I understand the rule better?}
\vocentry{894}{Невозможно}{nevaz'mozhnə}{It is impossible}{Невозможно быть успешным, не прилагая никаких усилий.}{It is impossible to be successful without making any effort.}
\vocentry{895}{Зеркало}{'zerkalə}{A mirror}{Я купила это красивое зеркало на распродаже.}{I bought this beautiful mirror at a sale.}
\vocentry{896}{Погибнуть}{pa'gibnut’}{To be killed, to perish (to die of unnatural causes)}{Мы не верили, что он мог погибнуть в автокатастрофе.}{We didn’t believe that he could be killed in a car crash.}
\vocentry{897}{Американский}{ameri'kanskij}{American}{У неё сильный американский акцент.}{She’s got a strong American accent.}
\vocentry{898}{Дым}{dym}{Smoke}{Я терпеть не могу сигаретный дым.}{I can’t stand cigarette smoke.}
\vocentry{899}{Гореть}{ga'ret’}{To burn}{Эта большая свеча может гореть до пяти часов.}{This big candle can burn up to five hours.}
\vocentry{900}{Плакать}{'plakat’}{To cry, to weep}{Я не могу не плакать когда смотрю эту сцену из фильма.}{I can’t help but cry when watching this movie scene.}
\vocentry{901}{Весьма}{ves’'ma}{Rather, quite}{Мой дедушка был весьма хорошо образован для своего времени.}{My grandfather was rather well-educated for his time.}
\vocentry{902}{Факт}{fakt}{A fact}{Тот факт, что я не замужем, меня совсем не беспокоит.}{The fact that I’m not married doesn’t bother me at all.}
\vocentry{903}{Двигаться}{'dvigatsa}{To move}{Как ты заставляешь эти картинки двигаться?}{How do you make these pictures move?}
\vocentry{904}{Рыба}{'ryba}{Fish}{Эта рыба не водится в солёной воде.}{This fish is not found in salty water.}
\vocentry{905}{Добавить}{da'bavit’}{To add}{Попробуй добавить в суп немного больше перца.}{Try to add a bit more pepper to the soup.}
\vocentry{906}{Удивиться}{udi'vitsa}{To be surprised, to wonder at}{Ты можешь удивиться, но я уже всё знаю.}{You may be surprised but I already know everything.}
\vocentry{907}{Бабушка}{'babushka}{A grandmother}{Моя бабушка готовит лучше любого шеф-повара.}{My grandmother cooks better than any chef.}
\vocentry{908}{Вино}{vi'no}{Wine}{Я предпочитаю полусухое вино.}{I prefer semidry wine.}
\vocentry{909}{Ибо}{'ibə}{For, since, as (literary, has an ironical connotation if used in everyday speech)}{Мы не остановимся, ибо наша борьба священна.}{We shall not stop for our struggle is sacred.}
\vocentry{910}{Учитель}{u'tchitel’}{A teacher}{Хороший учитель должен быть творческим человеком.}{A good teacher must be a creative person.}
\vocentry{911}{Действовать}{'dejstvavat’}{To act}{Врачи должны были действовать быстро, чтобы спасти ему жизнь.}{The doctors had to act fast to save his life.}
\vocentry{912}{Осторожно}{asta'rozhnə}{Carefully}{Она взяла кристальный шар и осторожно положила его в коробку.}{She took the crystal ball and put it into the box carefully.}
\vocentry{913}{Круг}{krug}{A circle; a range}{Круг – символ бесконечности.}{The circle is a symbol of infinity.}
\vocentry{914}{Папа}{'papa}{Dad}{Наш папа часто брал нас с собой на рыбалку.}{Our dad often took us fishing with him.}
\vocentry{915}{Правильно}{'pravil’nə}{Correctly}{Он ответил правильно на все вопросы и сдал тест.}{He answered all the questions correctly and passed the test.}
\vocentry{916}{Недавно}{ne'davnə}{Recently, lately}{Я видела его недавно, и он ничуть не изменился.}{I’ve seen him recently and he hasn’t changed a bit.}
\vocentry{917}{Держаться}{der'zhatsa}{To hold on, to hold on to}{Я знаю, нам всем тяжело, но мы должны держаться.}{I know it’s hard for all of us but we must hold on.}
\vocentry{918}{Причём}{pri'tchjom}{Moreover, with}{Эти близнецы окончили один университет, причём оба работают в одном месте.}{These twins graduated from the same university, moreover both work at the same place.}
\vocentry{919}{Лететь}{le'tet’}{To fly}{Крыло птицы было сломано, и она не могла лететь.}{The bird’s wing was broken and it couldn’t fly.}
\vocentry{920}{Носить}{na'sit’}{To wear; to carry}{Я всегда стараюсь носить шляпу в жаркую погоду.}{I always try to wear a hat in hot weather.}
\vocentry{921}{Повод}{'povəd}{A reason, an excuse}{Тебе не нужно искать повод, чтобы позвонить мне.}{You needn’t find a reason to call me.}
\vocentry{922}{Лагерь}{'lager’}{A camp}{Она поехала в летний лагерь на каникулах.}{She went to a summer camp on holidays.}
\vocentry{923}{Птица}{'ptitsa}{A bird}{Попугай – единственная птица, которая может говорить.}{The parrot is the only bird that can talk.}
\vocentry{924}{Корабль}{ka'rabl’}{A ship}{Этот старый корабль в основном используется для экскурсий вдоль берега.}{This old ship is mostly used for excursions along the shore.}
\vocentry{925}{Мнение}{'mnenije}{An opinion}{Твоё мнение много значит для меня.}{Your opinion means a lot to me.}
\vocentry{926}{Ночной}{natcth'noj}{Night, nightime (adj)}{Твой ночной звонок напугал меня.}{Your night call frightened me.}
\vocentry{927}{Здоровый}{zda'rovyj}{Healthy}{Овсянка с молоком – это здоровый завтрак.}{Oats with milk is a healthy breakfast.}
\vocentry{928}{Зима}{zi'ma}{Winter}{Зима – время семейных праздников и волшебства.}{Winter is the time of family holidays and magic.}
\vocentry{929}{Сухой}{su'hoj}{Dry}{Сухой климат хорошо влияет на моё здоровье.}{Dry climates have a good influence on my health.}
\vocentry{930}{Километр}{kila'metr}{A kilometer}{До заправки всего один километр.}{There’s only one kilometer to the filling station.}
\vocentry{931}{Кровать}{kra'vat’}{A bed}{Наша новая кровать более просторная.}{Our new bed is more spacious.}
\vocentry{932}{Привыкнуть}{pri'vyknut’}{To get used to}{Я не могу привыкнуть делать зарядку по утрам.}{I can’t get used to doing morning exercises.}
\vocentry{933}{Прочее}{'protcheje}{So on, so forth}{Здесь я храню всякие мелочи: серёжки, кольца и прочее.}{Here I keep different little things: earrings, rings, and so on.}
\vocentry{934}{Свободный}{sva'bodnyj}{Free; vacant}{Свободный человек не зависит от мнения других людей.}{A free person doesn’t depend on the opinion of other people.}
\vocentry{935}{Лестница}{'lesnitsa}{Stairs; a ladder}{Эта лестница старая и всё время скрипит.}{These stairs are old and are creaking all the time.}
\vocentry{936}{Неужели}{neu'zheli}{Really (to express surprise)}{Мама, неужели эта девочка на фото – это ты?}{Mom, is this girl in the photo really you?}
\vocentry{937}{Обязательно}{abja'zatel’nə}{Necessarily}{Иметь образование не обязательно значит быть умным.}{Having an education doesn’t necessarily mean being clever.}
\vocentry{938}{Вверх}{vverh}{Up, upwards}{Мы шли вверх по холму более двух часов и очень устали.}{We’d been walking up the hill for more than two hours and were very tired.}
\vocentry{939}{Детство}{'detstvə}{Childhood}{Моё детство ассоциируется у меня с теплом и радостью.}{I associate my childhood with warmth and joy.}
\vocentry{940}{Остров}{'ostrəv}{An island}{Когда я был маленьким мальчиком, я мечтал попасть на необитаемый остров.}{When I was a little boy I dreamt of getting to an uninhabited island.}
\vocentry{941}{Статья}{stat’'ja}{An article}{Эта статья о детском доме сильно впечатлила меня.}{That article about an orphanage impressed me a lot.}
\vocentry{942}{Позвонить}{pazva'nit’}{To call (by phone), (perfect)}{Не забудь позвонить мне, когда вернёшься домой.}{Don’t forget to call me when you return home.}
\vocentry{943}{Столь}{stol’}{So, such (literary style)}{Его цели были столь амбициозными, что никто не верил в его успех.}{His goals were so ambitious that nobody believed in his success.}
\vocentry{944}{Мешать}{me'shat’}{To bother, to prevent; to stir}{Дети, постарайтесь не мешать мне, пока я работаю.}{Kids, try not to to bother me while I’m working.}
\vocentry{945}{Водка}{'vodka}{Vodka}{Водка – это крепкий алкогольный напиток.}{Vodka is a strong alcoholic beverage.}
\vocentry{946}{Темнота}{temna'ta}{Darkness}{Темнота пугает многих людей.}{Darkness frightens many people.}
\vocentry{947}{Возникнуть}{vaz'niknut’}{To arise, to emerge (perfective)}{При решении этой задачи могут возникнуть различные вопросы.}{While solving this task, different questions may arise.}
\vocentry{948}{Способный}{spa'sobnyj}{Capable, able}{Нам нужен человек, способный решить эту проблему самостоятельно.}{We need a person who’s capable of solving this problem independently.}
\vocentry{949}{Станция}{'stantsija}{A station}{Эта железнодорожная станция закрыта на ремонт.}{This railway station is closed for repairs.}
\vocentry{950}{Желать}{zhe'lat’}{To wish, to desire}{У меня есть всё. Я не могу желать большего.}{I’ve got everything. I can’t wish for more.}
\vocentry{951}{Попробовать}{pap'robavat’}{To try; to taste (to test the flavor of something) (perfective)}{Я хочу попробовать выучить ещё один иностранный язык.}{I want to try to learn another foreign language.}
\vocentry{952}{Получаться}{palu'tchatsa}{To be able to, to work out well}{Со временем это упражнение будет получаться у тебя всё лучше и лучше.}{With time you will be able to do this exercise better and better.}
\vocentry{953}{Гражданин}{grəzhda'nin}{A citizen}{Каждый гражданин имеет право на отдых.}{Every citizen has the right to rest.}
\vocentry{954}{Странно}{'strannə}{Weirdly, strangely}{Что случилось? Ты ведёшь себя странно.}{What’s happened? You behave weirdly.}
\vocentry{955}{Вскоре}{'vskore}{Soon (a more literary variant of ‘скоро’)}{Вскоре они совсем забыли о своей ссоре.}{Soon they forgot about their argument.}
\vocentry{956}{Команда}{ka'manda}{A team, a crew; an order}{Наша футбольная команда проиграла в финальном матче.}{Our football team lost in the final match.}
\vocentry{957}{Заболевание}{zabale'vanije}{A disease, a sickness}{Учёные научились контролировать это заболевание.}{Scientists learned to control this disease.}
\vocentry{958}{Живот}{zhi'vot}{Belly, stomach}{Мой толстый живот меня раздражает! Я хочу похудеть.}{My fat belly irritates me! I want to lose weight.}
\vocentry{959}{Ставить}{s'tavit’}{To put, to place}{Я часто забываю ставить вещи обратно на свои места.}{I often forget to put things back to their places.}
\vocentry{960}{Ради}{'radi}{For the sake of}{Они решили не разводиться ради детей.}{They decided not to get divorced for the sake of the children.}
\vocentry{961}{Тишина}{tishi'na}{Silence, quietness}{Мне нужна полная тишина, чтобы сосредоточиться.}{I need complete silence to concentrate.}
\vocentry{962}{Понятно}{pa'njatnə}{Clearly}{Ты не умеешь понятно объяснять – я ничего не понимаю.}{You aren’t able to explain clearly – I don’t understand anything.}
\vocentry{963}{Фронт}{front}{A front (military)}{Во время Второй мировой войны даже подростки уходили на фронт.}{During World War II even teenagers went to the front.}
\vocentry{964}{Щека}{stche'ka}{A cheek}{После укуса её щека некоторое время оставалась красной.}{After the sting her cheek remained red for some time.}
\vocentry{965}{Страшно}{'strashnə}{One is scared (with the pronoun in dative case); scary}{Мне страшно! Включи свет!}{I am scared! Turn on the light!}
\vocentry{966}{Район}{ra'jon}{A district; an area}{Этот район самый густонаселённый в городе.}{This district is the most densely populated one in the city.}
\vocentry{967}{Наверно}{na'vernə}{Probably, likely}{Он, наверно, проспал. Вчера он очень поздно вернулся домой.}{He probably overslept. He came back home very late yesterday.}
\vocentry{968}{Проводить}{prava'dit’}{To spend, to pass}{Я люблю проводить время с родственниками.}{I like to spend time with my relatives.}
\vocentry{969}{Выражение}{vyra'zhenije}{An expression}{Выражение её лица объясняло всё без слов.}{The expression on her face explained everything without words.}
\vocentry{970}{Слегка}{sleh'ka}{Slightly, a bit}{Наши планы на лето слегка изменились.}{Our plans for the summer have changed slightly.}
\vocentry{971}{Мешок}{me'shok}{A sack}{Он принёс целый мешок подарков.}{He brought a whole sack of presents.}
\vocentry{972}{Обещать}{obe'stchat’}{To promise}{Я не могу обещать тебе, что починю твою машину, но я попробую.}{I can’t promise you to repair your car but I’ll try.}
\vocentry{973}{Дорогой}{dara'goj}{Expensive; dear}{Этот кожаный кошелёк слишком дорогой для меня.}{This leather purse is too expensive for me.}
\vocentry{974}{Судить}{su'dit’}{To judge}{Я стараюсь не судить людей по внешности.}{I try not to judge people by appearance.}
\vocentry{975}{Большинство}{bal’shinst'vo}{Majority, most}{Большинство студентов приняли участие в конкурсе.}{The majority of the students took part in the contest.}
\vocentry{976}{Собраться}{sab'ratsa}{To gather, to assemble (perfective)}{Мы решили собраться после работы и обсудить планы на выходные.}{We decided to gather after work and to discuss the plans for the weekend.}
\vocentry{977}{Управление}{uprav'lenije}{Administration, managing}{Эффективное управление любой работой влияет на результат.}{Efficient administration of any work influences the result.}
\vocentry{978}{Колоть}{ka'lot’}{To chop (firewood); to prick}{Я не умею колоть дрова.}{I can’t chop firewood.}
\vocentry{979}{Мокрый}{'mokryj}{Wet, damp}{Ты оставил ковёр под дождём, и теперь он мокрый насквозь.}{You’ve left the carpet under the rain and now it’s wet through.}
\vocentry{980}{Приказ}{pri'kaz}{An order, a command}{Солдаты получили приказ наступать.}{The soldiers got the order to advance.}
\vocentry{981}{Прямой}{prja'moj}{Straight}{У него был длинный прямой нос и тонкие губы.}{He had a long straight nose and thin lips.}
\vocentry{982}{Закричать}{zakri'chat’}{To cry out (perfective)}{Она хотела закричать от радости, но вспомнила что все спят.}{She wanted to cry out with joy but remembered that everyone was sleeping.}
\vocentry{983}{Кончиться}{'kontchitsa}{To end; to run out of (perfective)}{Это приключение просто не могло кончиться хорошо.}{That adventure just couldn’t end well.}
\vocentry{984}{Куст}{kust}{A bush}{Мне нравится, как цветёт этот куст у моего окна.}{I like the way this bush by my window blooms.}
\vocentry{985}{Стрелять}{stre'ljat’}{To shoot, to fire}{Я не могу стрелять в животных.}{I can’t shoot animals.}
\vocentry{986}{Художник}{hu'dozhnik}{An artist, a painter}{Этот художник скоро открывает новую выставку.}{This artist is opening a new exhibition soon.}
\vocentry{987}{Знак}{znak}{A sign, a mark}{Ты знаешь, для чего используется этот знак?}{Do you know what this sign is used for?}
\vocentry{988}{Завод}{za'vod}{A factory, a plant}{Этот завод загрязняет окружающую среду.}{This factory contaminates the environment.}
\vocentry{989}{Кулак}{ku'lak}{A fist}{Мальчик сжал кулак и погрозил обидчику.}{The boy clenched his fist and shook it at the offender.}
\vocentry{990}{Использовать}{is'pol’zavat’}{To use}{Мы должны эффективно использовать эту новую возможность.}{We must efficiently use this new opportunity}
\vocentry{991}{Стакан}{sta'kan}{A glass}{Твой стакан наполовину пуст или наполовину полон?}{Is your glass half empty or half full?}
\vocentry{992}{Пахнуть}{'pahnut’}{To smell (of)}{Что может так хорошо пахнуть? У меня слюнки текут!}{What can smell so good? My mouth is watering!}
\vocentry{993}{Отсюда}{ats'juda'}{Out of here, from here}{Выведи собаку отсюда! Она вся грязная после прогулки.}{Get the dog out of here! It’s all dirty after the walk.}
\vocentry{994}{Рот}{rot}{A mouth}{Она открыла рот, чтобы возразить, но передумала.}{She opened her mouth to contradict but changed her mind.}
\vocentry{995}{Пора}{pa'ra}{It is time; a season}{Тебе пора найти работу и начать жить самостоятельно.}{It is time for you to find a job and to start living independently.}
\vocentry{996}{Передать}{pere'dat’}{To hand over, to pass}{Я уже стар, и решил передать бизнес своим детям.}{I’m old and have decided to hand the business over to my children.}
\vocentry{997}{Лечение}{le'tchenije}{Treatment}{Это лечение очень дорогое, но эффективное.}{This treatment is very expensive but very efficient.}
\vocentry{998}{Высший}{'vysshij}{The highest}{Она получила высший балл за экзамен по математике.}{She got the highest grade for her Math exam.}
\vocentry{999}{Сутки}{'sutki}{24 hours, day and night}{Она заблудилась в лесу и нашла выход только через сутки.}{She got lost in the woods and found the way out only in 24 hours.}
\vocentry{1000}{Вставать}{vsta'vat’}{To stand up, to get up}{Инструктор предупредил нас, что вставать в лодке очень опасно.}{The instructor warned us that it’s very dangerous to stand up in a boat.}
\end{multicols}
